\chapter{Differential Propagation and Survival}
A differential leaving $h_6^2\in E_2^{2,128}(S^0)$ has its target in
stem $125$. The calculations in this chapter constrain that stem and the
adjacent stems needed for the extension argument. The central classes
are $W,T,U$ and $H$ of Chapter~3: $T$ is the possible last target,
$W$ is the possible source of a $d_6$ hitting $T$, and $U$ enters the
theta-five square calculation. We also need the bottom-cell images
in $C\nu$ and $C_2$.

We begin with the finite differential assumptions, then calculate the
surviving groups and the possible incoming differentials to $T$.
Fact~\ref{paper:7.6} summarizes the sphere information;
Theorem~\ref{sem:Computation:cofiber-results} gives the cofiber
differential and excludes incoming differentials to $i_*T$ through
page five.

The calculations use the complete $E_2$ bases of Chapter~3. A basis
expression always denotes a second-page label; an equation $d_r(x)=y$
asserts that both labels have representatives on page $r$ satisfying
that equation.
As in Chapter~1, a label \emph{reaches} a page if it has a
representative there, which may be zero. \emph{Nonzero survival}
requires a nonzero representative. A permanent label has no outgoing
differential, although an incoming boundary may make its limiting class
zero. The later permanence argument excludes subsequent outgoing differentials;
nonzero survival also uses incoming-differential control.

The sphere and Hopf-cofiber inputs are stated first, followed by their
consequences for the named classes. References for the calculations
are collected in the final section.

\section{Differential assumptions and page bounds}


\begin{theorem}[Finite differential assumptions]
\label{sem:SSComputation:admitted-finite-inputs}
\label{sem:Computation:table-certification}
\label{sem:SSComputation:ordinary-differential-inputs}
\level{coarse}
\uses{sem:AdamsComputation:moore-object,sem:Computation:table-bases,sem:Computation:w-t-u,sem:SpectralSequence:differential,sem:SpectralSequence:cycles-boundaries,sem:SpectralSequence:representative-expansion}
Assume the $200$ statements of the
\href{paper/finite-calculations.html}{supplementary finite calculation table}.
All $200$ are distinct coordinate rows. They include the actual bottom-cell
class of $C_2$ and the fixed sphere class $V$:
\[
 d_4(\bar x)=0,\quad \bar x\in E_2^{8,132}(C_2),\qquad
 d_7(V)=0,\quad V\in E_2^{9,132}(S^0).
\]
Each equation includes representatives of the specified labels on its page.
The first equation has target degree $(12,135)$ and the second $(16,138)$.
The Lean certification of these $200$ computational inputs is pending.
As in Chapter~3, indices
begin at $0$, a list denotes the sum of the corresponding basis
vectors, and the empty list denotes zero. For $C_2$, coordinates denote
images of sphere basis vectors under the bottom-cell map $i_*$. The last
two rows have source coordinates $[0]$ on $C_2$ and $[0,1]$ on the sphere.
Each statement is a
differential equation, continuation to a page $r\ge2$, a boundary bound,
or the exclusion of a differential equation. A differential with source
bidegree $(s,t)$ has target bidegree $(s+r,t+r-1)$. The table includes
page-$1000$ continuations.
A boundary bound records membership in the indicated boundary group,
allowing zero and sums of earlier boundaries. An excluded equation is
the negation of the specified representative differential relation.

The four direct $d_2$ equations of
Theorem~\ref{sem:AdamsComputation:direct-d2-records} and the two staircase
assumptions in Theorem~\ref{sem:SSComputation:staircase-inputs} are separate
inputs. When a propagation argument needs one of these four second
differentials, it uses the Chapter~3 input directly; none is repeated in
the finite table above.
\lean{KIP126.SSComputation.Internal.Certification.certified_result,KIP126.Computation.Table.result,KIP126.SSComputation.Moore.bottom_source_d4_zero,KIP126.SSComputation.Moore.v_d7_zero}
\end{theorem}

\begin{theorem}[Finite-page survival of the selected classes]
\label{sem:SSComputation:selected-class-survival}
\level{coarse}
\uses{sem:SSComputation:ordinary-differential-inputs,sem:AdamsComputation:moore-object,sem:Computation:table-bases,sem:SpectralSequence:cycles-boundaries}
The two ordinary equations in the finite input table imply
\[
 d_3(\bar x)=0,\qquad \bar x\in Z_4^{8,132}(C_2),\qquad
 V\in Z_8^{9,132}(S^0).
\]
Thus $\bar x$ reaches $E_5$ and $V$ reaches $E_9$.
\lean{KIP126.SSComputation.Moore.bottom_source_d3_zero,KIP126.SSComputation.Moore.bottom_source_cycles_four,KIP126.SSComputation.Moore.v_cycles_eight}
\end{theorem}
\begin{proof}
The fourth-page representatives for $\bar x$ give the preceding zero
third differential, and $d_4(\bar x)=0$ gives the fifth-page lift.
The zero seventh differential gives an eighth-page representative for $V$.
Both basis classes at its next possible target $(17,139)$ are killed by
the two third differentials from $(14,137)$ in the finite table.
The eighth differential is therefore zero and $V$ reaches $E_9$.
\end{proof}

\begin{theorem}[Representatives for two staircase classes]
\label{sem:SSComputation:staircase-inputs}
\level{coarse}
\uses{sem:Computation:table-bases,sem:SpectralSequence:cycles-boundaries}
Assume that the two sphere labels
\[
 b_{S^0;14,138,2},\qquad
 b_{S^0;15,139,1}+b_{S^0;15,139,3}
\]
reach page $1000$. Equivalently, they belong respectively to
$Z_{999}^{14,138}(S^0)$ and $Z_{999}^{15,139}(S^0)$.
These are two additional assumptions beyond the $200$ statements above.
The page-$1000$ representatives are allowed to be zero.
\lean{KIP126.Computation.Table.staircase_result}
\end{theorem}
The passage from these finite page bounds to permanence uses the
vanishing of every later differential target. Nonzero survival also
requires control of incoming differentials.


\section{Sphere Adams differentials}


\begin{theorem}[Third-page sphere differentials]
\label{sem:Computation:low-differentials}
\level{coarse}
\uses{sem:Computation:standard-atoms,sem:Computation:named-basis-identities}
The sphere Adams spectral sequence has the following nonzero
differentials:
\[
d_3(h_4x_{109,12})=h_1x_{122,15,2},\qquad
 d_3(h_0^2x_{123,13,2})=h_0^2x_{122,16},
\]
\[
d_3(x_{126,4})=h_0^2x_{125,5}.
\]
\lean{KIP126.Computation.Standard.d3_h4_x109_12,KIP126.Computation.Standard.d3_h0Square_x123_13_2,KIP126.Computation.Standard.d3_x126_4}
\end{theorem}
\begin{proof}
\uses{sem:Computation:table-bases,sem:Computation:table-certification,sem:SSComputation:staircase-inputs}
In the bases of Chapter~3, these are the corresponding third
differentials in the assumed finite calculation. The basis identities
identify their sources and targets with the displayed products; the
targets are nonzero on $E_3$. Thus each source direction disappears
on the next page, while each target becomes a boundary. The second
differentials on $h_0^6h_6$ and $h_6$ are supplied independently in
Chapter~3.
\end{proof}

\begin{theorem}[A third-differential alternative]
\label{sem:Computation:d3-candidates}
\level{coarse}
\uses{sem:Computation:w-t-u,sem:Computation:named-basis-identities,sem:Computation:table-certification}
The differential on $x_{126,6}$ has one of the two targets
\[
 d_3(x_{126,6})=B=h_5x_{94,8}
 \quad\text{or}\quad
 d_3(x_{126,6})=B'=h_5x_{94,8}+h_6b.
\]
\lean{KIP126.Computation.Standard.d3_x126_6_candidates}
\end{theorem}

\section{Survival, incoming control and higher-page vanishing}

\begin{theorem}[The correction term survives nontrivially]
\label{sem:Computation:correction-survival}
\level{coarse}
\uses{sem:Computation:w-t-u,sem:SpectralSequence:permanent,sem:Computation:table-certification}
The class $C=e_0\Delta h_6g$ survives to a nonzero element of
$E_\infty^{13,137}$. In particular it provides the filtration-thirteen
correction class in the theta-five square calculation.
\lean{KIP126.Computation.Standard.correction_permanent}
\end{theorem}

% AimPaper: 7.6 fact 2156 2167
\begin{fact}[Sphere classes and the incoming differentials to $T$]
\label{paper:7.6}
\label{fact:theta5sqAF}
\label{sem:Computation:incoming-control}
\source{AimPaper Fact 7.6}
\level{coarse}
\uses{sem:Computation:table-bases,sem:Computation:table-certification,sem:Computation:w-t-u,sem:AdamsSynthetic:standard-classes,sem:Computation:square-line,sem:AdamsComputation:direct-d2-records}
The sphere computations of \cite{LWX126}, Fact 7.6, give:
\begin{enumerate}
\item $W=x_{126,8,4}+x_{126,8}$ survives nontrivially to $E_6$.
\item $T=h_1h_4x_{109,12}$ is a permanent label, with no outgoing
 differential. Its only possible nonzero incoming differentials are
\[
 d_6(W)=T\quad\text{or}\quad d_{12}(h_6^2)=T.
\]
Consequently no other page $r\ge2$ hits $T$, no differential from
$x_{126,6}$ has target $T$, and the $d_6$ target of $W$ belongs to
$\{0,T\}$.
\item $U=h_0^2x_{124,8}$ survives nontrivially to $E_\infty$.
\item $E_5^{25,150}(S^0)$ has a unique nonzero class, represented by
$H=g^4\Delta h_1g$.
\end{enumerate}
\lean{KIP126.Computation.Standard.w_to_e6,KIP126.Computation.Standard.t_permanent_cycle,KIP126.Computation.Standard.t_only_incoming,KIP126.Computation.Standard.u_permanent,KIP126.Computation.Standard.high_e5,KIP126.Computation.Standard.no_other_incoming,KIP126.Computation.Standard.table9,KIP126.Computation.Standard.w_d6_targets}
\end{fact}
\begin{proof}
\textit{(See \cite{LWX126}, Fact 7.6.)}
\href{paper/main.tex.html#L2156}{Original statement and context}.
An incoming differential to $T\in E_r^{14,139}$ has source degree
$(14-r,140-r)$. Figure 1 and the Appendix tables of \cite{LWX126}, in
stems $124$, $125$ and $126$, leave only the sources of filtrations eight and two, giving
the possible $d_6$ and $d_{12}$ above. The filtration-six source in
stem $126$ supports the nonzero $d_3$ of
Theorem~\ref{sem:Computation:d3-candidates}, so it cannot give another
incoming differential. The page-twelve source lies in degree $(2,128)$; by
Theorem~\ref{sem:Computation:square-line}, its only nonzero label is
$h_6^2$. The same calculation gives $W$ nonzero on $E_6$ and no outgoing
differential on $T$.

In stem $124$, the class $U$ has a nonzero limiting image. In stem $125$
and filtration twenty-five, quotienting by the earlier differentials
leaves a one-dimensional fifth-page group generated by $H$. These
calculations use the differential assumptions and the complete bases,
independently of the permanence of $h_6^2$.
\end{proof}

\begin{theorem}[Boundary products of the permanent-class candidates]
\label{sem:Computation:product-boundaries}
\level{coarse}
\uses{sem:Computation:w-t-u,sem:Computation:standard-atoms,sem:SpectralSequence:cycles-boundaries}
The products $Ph_2$ and $Qh_2$ are $d_2$ boundaries, and hence vanish
on $E_3$.
\lean{KIP126.Computation.Standard.p_h2_boundary,KIP126.Computation.Standard.q_h2_boundary}
\end{theorem}
\begin{proof}
\uses{sem:AdamsComputation:direct-d2-records,sem:Computation:named-basis-identities,sem:SpectralSequence:cycles-boundaries}
The second $d_2$ equation of Chapter~3 has target
$b_{S^0;12,137,2}=Ph_2$. Adding the third and fourth equations gives
target $b_{S^0;13,138,3}+b_{S^0;13,138,4}=Qh_2$, by the product
identities in that chapter. Thus both products lie in the incoming
$d_2$ image. Their vanishing on $E_3$ will constrain the products of
later synthetic lifts.
\end{proof}


\begin{theorem}[The required higher-page vanishing components]
\label{sem:Computation:higher-page-vanishing}
\level{coarse}
\uses{sem:AdamsSynthetic:families,sem:Computation:bounded-comparison}
The following groups vanish: $E_5^{11,135}$, $E_4^{12,136}$, $E_4^{12,137}$ and
$E_5^{13,138}$.
\lean{KIP126.Computation.Standard.e5_stem124_af11,KIP126.Computation.Standard.e4_stem124_af12,KIP126.Computation.Standard.e4_stem125_af12,KIP126.Computation.Standard.e5_stem125_af13}
\end{theorem}
\begin{proof}
\uses{sem:Computation:table-bases,sem:Computation:table-certification,sem:AdamsComputation:direct-d2-records}
At each step $E_{r+1}=\ker d_r/\operatorname{im}d_r$. In these
bidegrees, the direct $d_2$ equations of Chapter~3 and the subsequent
finite differential calculations make the incoming images exhaust
the cycles by the stated page. The resulting quotients are zero.
The complete bases ensure that this calculation accounts for every
cycle direction.
\end{proof}


\section{The Hopf cofiber differential}

\begin{fact}[Nonzero survival of $Y$ to $E_5$]
\label{sem:Computation:y-to-e5}
\source{AimPaper Fact 7.15, survival assertion}
\level{coarse}
By \cite{LWX126}, Fact 7.15, the class
$Y=h_0^2x_{125,9,2}$ in bidegree $(11,136)$ of the sphere Adams
$E_2$-page survives nontrivially to $E_5$.
\uses{sem:Computation:w-t-u,sem:Computation:table-bases,sem:Computation:table-certification,sem:AdamsComputation:direct-d2-records}
\lean{KIP126.Computation.Standard.y_to_e5}
\end{fact}
\begin{proof}
The stem-$125$ calculation in the Appendix of \cite{LWX126} leaves
$Y$ nonzero on $E_5$. The exclusion of all incoming differentials uses
Proposition~\ref{paper:7.8}, as explained in Chapter~6.
\end{proof}

\begin{theorem}[The cofiber differential and survival window]
\label{sem:Computation:cofiber-results}
\level{coarse}
\uses{sem:Computation:nu-cofiber,sem:Computation:w-t-u}
In the Hopf cofiber Adams spectral sequence there is a nonzero
differential $d_3(A_{C\nu})=i_*Y$ in bidegrees
$(8,134)\to(11,136)$. The class $i_*T$ survives nontrivially to page $6$
and is not hit on any page in the window $2\leq r\leq5$.
\lean{KIP126.Computation.Standard.cofiber_d3,KIP126.Computation.Standard.cofiber_target_to_e6,KIP126.Computation.Standard.cofiber_target_not_hit_through5}
\end{theorem}
\begin{proof}
\uses{sem:Literature:nu-cofiber,sem:Computation:table-bases,sem:Computation:named-basis-identities,sem:Computation:table-certification,sem:Computation:y-to-e5}
The basis identities of Chapter~3 are
\[
 A_{C\nu}=b_{C\nu;8,134,0}+b_{C\nu;8,134,3}+b_{C\nu;8,134,4},
 \qquad i_*Y=b_{C\nu;11,136,3}.
\]
In these bases, the cofiber differential sends that sum of three source
vectors to the indicated target vector. This is the asserted equation.

The exact sequence of the Hopf cofiber has the segment
\[
 E_2^{10,132}(S^0)\xrightarrow{\ h_2\ }
 E_2^{11,136}(S^0)\xrightarrow{\ i_*\ }E_2^{11,136}(C\nu).
\]
The first group is zero, so $i_*$ is injective.
Since $Y$ is nonzero, so is $i_*Y$. All five basis elements of
$E_2^{9,135}(C\nu)$ reach $E_3$, so their $d_2$ images vanish.
There are therefore no $d_2$ boundaries in the target, and every
third-page representative of $i_*Y$ is nonzero.

Thus the target of the displayed third differential is nonzero on
its page. The cofiber calculation also leaves $i_*T$ nonzero through
page six and excludes incoming differentials to it on pages two
through five;
the bottom-cell comparison identifies this class with the image of $T$.
\end{proof}



\section{Third-page screening of the second line}


\begin{theorem}[The second line after the second differential]
\label{sem:AdamsCalculations:two-line-d2-screening}
\level{coarse}
The second-line calculation uses \cite{Adams60}, Theorem 4.6.1,
and \cite{LinExt4}, Theorem 1.3; compare the proof of Corollary 1.5
in \cite{LWX126}.

Write $H_{i,j}=h_ih_j$. If $0\leq i\leq j\leq7$,
$H_{i,j}\ne0$ on $E_2$, and $H_{i,j}$ reaches $E_3$, then its label
belongs to the following list:
\[
 H_{0,2},H_{0,3},H_{2,4},H_{1,j}\ (3\leq j\leq7),h_j^2\ (0\leq j\leq7),
 H_{2,5},H_{3,5},H_{3,6},H_{0,4}.
\]
These are the second-line $d_2$ candidates reaching $E_3$.
\lean{KIP126.SSComputation.two_line_d2_candidates}
\end{theorem}
\begin{proof}
\uses{sem:Literature:one-line-differentials,sem:Computation:bounded-comparison,sem:Computation:table-bases,sem:Computation:named-basis-identities}
There are 36 pairs in the stated range. Twenty already belong
to the candidate list. The seven adjacent products $H_{i,i+1}$ vanish
because their second-page groups are zero, contradicting the
nonzero hypothesis.

For the remaining nine pairs, multiplicative Leibniz gives
\[
\begin{aligned}
d_2(H_{0,j})&=h_0^2h_{j-1}^2 &&(j=5,6,7),\\
d_2(H_{2,j})&=h_0h_2h_{j-1}^2 &&(j=6,7),\\
d_2(H_{3,7})&=h_0h_3h_6^2,\\
d_2(H_{4,6})&=h_0h_3^2h_6,\\
d_2(H_{4,7})&=h_0h_3^2h_7,\\
d_2(H_{5,7})&=h_0h_4^2h_7.
\end{aligned}
\]
The other Leibniz terms vanish by adjacent-product relations and
$h_4h_6^2=0$. Each displayed target is a nonzero basis
monomial under the multiplicative $E_2$ comparison of Chapter~3.
A label that reaches the third page has zero second differential,
so these nine pairs are excluded. The list therefore retains $H_{0,4}$
and all squares in the stated range. Later differentials and the
geometric exclusion impose further restrictions.
\end{proof}


\section{Selected classes, survival ranges and incoming differentials}

% AimPaper: 7.13 fact 2377 2385
\begin{fact}[The second differential on $x_{125,8}$]
\label{paper:7.13}
\label{fact:x_123_9}\label{fact:x1239}
\source{AimPaper Fact 7.13, second assertion}
\level{coarse}
The sphere Adams spectral sequence has the nonzero differential
\[
 d_2(x_{125,8})=h_1V+U.
\]
The continuation of $V$ to $E_9$ used below is supplied by
Theorem~\ref{sem:SSComputation:selected-class-survival}.
\uses{sem:Computation:table-bases,sem:Computation:table-certification}
\lean{KIP126.Computation.Standard.d2_x125_8}
\end{fact}

% AimPaper: 7.19 fact 2613 2615
\begin{fact}[The class X through page six]
\label{paper:7.19}
\source{AimPaper Fact 7.19}
\level{coarse}
From \cite{LWX126}, Fact 7.19.\par

\label{fact:h_1_mul_x_121_7}\label{fact:h1x1217}
    The class $X=h_1x_{121,7}$ survives nontrivially to $E_6$.
    No classical incoming differential hits this label.

\uses{sem:Computation:table-bases,sem:Computation:table-certification}


\lean{KIP126.Computation.Standard.x_to_e6,KIP126.Computation.Standard.x_not_hit}

\end{fact}
\begin{proof}
\textit{(See \cite{LWX126}, Fact 7.19.)} \href{paper/main.tex.html#L2613}{Original statement and context}.

The stem-$122$ calculation in the Appendix of \cite{LWX126} gives
no outgoing differential on $X$ before $d_6$ and no incoming
differential to $X$. Thus $X$ is nonzero on $E_6$. These sphere
calculations will be used in constructing the later $\nu$-extension.
\end{proof}

\section{Permanent cycles in stem 122}

% AimPaper: 7.21 fact 2689 2696
\begin{fact}[Nonzero permanent classes in stem 122]
\label{paper:7.21}
\source{AimPaper Fact 7.21}
\level{coarse}
From \cite{LWX126}, Fact 7.21.\par

\label{fact:stem122}
\begin{enumerate}
    \item $P=h_6Md_0$ survives nontrivially to $E_\infty$.
    \item $Q=h_5x_{91,11}$ survives nontrivially to $E_\infty$.
\end{enumerate}

\uses{sem:Computation:table-bases,sem:Computation:table-certification}


\lean{KIP126.Computation.Standard.p_permanent,KIP126.Computation.Standard.q_permanent}

\end{fact}
\begin{proof}
\textit{(See \cite{LWX126}, Fact 7.21.)} \href{paper/main.tex.html#L2689}{Original statement and context}.

The stem-$122$ calculation in the Appendix of \cite{LWX126} leaves
$P$ and $Q$ as nonzero permanent classes. Their products with $h_2$
are the second-page boundaries computed above. The coexistence of
these permanent sources with those boundary products is the sphere
information used in the contradiction argument; the passage to their
synthetic representatives uses Chapter~4.
\end{proof}

\section{Sources for the finite calculations}

The propagation calculations of \cite{LinProgram} use products,
naturality and cofiber sequences; see also
\cite{LWXZenodoCurrent,LWXMachine}. They supply the finite inputs of
this chapter, separately from the four direct second-differential
inputs of Chapter~3.

The \href{paper/finite-calculations.html}{finite calculation table}
contains the $200$ distinct assumptions used in this chapter, expressed
in the complete bases of Chapter~3.

The two staircase labels of
Theorem~\ref{sem:SSComputation:staircase-inputs} occur in the sphere
calculation of \cite{LWXZenodoCurrent}, version
\texttt{v126.3.cw49}. The permanence convention of
\cite{LWXMachine}, Notation 3.5, assumes that visible nonzero
differentials have length less than $1000$. The additional assumptions
used here are precisely that these two labels have representatives on
page $1000$. The later permanence argument excludes every subsequent differential target.
